% !TeX root = ../../Book1.tex
% 纲要标题：## 附录 A　公理化几何与几何史的方法
% 纲要位置：几何学计划纲要.md，第 920 行。
% 正文状态：初稿；尚未正式审核。

% 纲要内容开始
% 
% > 选读；保留公理、模型、坐标化之间的逻辑辨析，不充当正文证明的替代。
% 
% * **公理与模型**
%   * 关联、顺序、合同及连续性
%   * 平行公设和非欧模型
%   * 独立性与相对一致性
% 
% * **坐标化及历史论证**
%   * 几何对象与代数模型的双向建立
%   * 历史证明的原始语境与现代重述
% 
% ---
% 
% 纲要内容结束

\chapter{公理化几何与几何史的方法}
\label{g1:app:A}
\BookChapterNumber{g1:app:A}

\begin{chapterabstract}
	几何证明既要说明从哪些假设出发，也要说明“点”“直线”“相等”等词在所用模型中指什么。本附录以有限组关联公理、欧氏与双曲平面、仿射坐标和等边三角形构造为例，辨别公理、模型与计算之间的关系。我们在明确列出的公理系统中比较平行命题的不同真值，解释模型如何支持独立性与相对一致性的判断；再把已有的仿射及欧氏结构写成坐标，并核对坐标变换保留什么。历史部分回到具体原始论证，区分原文的构造条件、现代补充的假设和另一种证明。
\end{chapterabstract}

\begin{learninggoals}
	\begin{itemize}
		\item 区分点线关联、介于关系、合同关系与连续性分别规定的内容，并在具体模型中核验所列条件。
		\item 在欧氏与双曲平面中判断给定点线之间的平行关系，说明模型中使用的直线与图纸上的形状有何区别。
		\item 用同一公理系统的两个模型解释独立性，区分模型存在、形式推导与依赖背景理论的相对一致性。
		\item 由实仿射标架建立坐标，恢复直线、介于及内积数据，核对换标架时对象和矩阵怎样变化。
		\item 阅读一个历史几何构造时，分别检查目标、已有命题、对象存在性和现代重述中增加的条件。
	\end{itemize}
\end{learninggoals}

在一条线段的两个端点各画一个圆，半径都取这条线段的长度；若两圆相交，交点与两个端点便组成等边三角形。最后一句的长度比较很直接，真正需要停下来说明的是前一句中的“相交”。图纸能够帮助找到候选位置，却不能替代所选几何体系中的存在性依据。若改用坐标求解，方程给出候选点后，仍须检查这些坐标属于允许使用的数系。一个熟悉的作图问题就这样同时涉及构造、代数和公理。

\ProseParagraphBreak

本卷正文主要在已经指定结构的空间中工作：实仿射空间给出点差和直线，内积给出距离，流形与度量给出局部坐标和测地线。公理化研究则允许先把点、线及其关系视为尚未解释的对象，问哪些假设足以推出一项结论。模型再为这些词指定具体含义，使每条假设成为能够核验的数学陈述。两种做法解决的问题不同，但可以相互校验：模型检查假设是否相容，也提醒我们一个证明究竟用了哪些额外结构。

\ProseParagraphBreak

这里使用\chapref{g1:ch:01}的实仿射空间、内积与标架，以及\secref{g1:sec:1001}已经构造并证明性质的双曲平面。逻辑部分只介绍当前模型论证所需的满足关系与推理保真性，不预设数理逻辑课程。坐标部分从已有实仿射结构出发，不把选取标架等同于从任意抽象点线系统重建一个数域。历史部分采用可定位的具体文本，依其所陈述的条件比较证明，而不从现代计算反推作者当时的想法。

\section{公理与模型}
\label{g1:sec:A01}

“两点确定一条直线”可以在纸上画出的直线中理解，也可以在双曲平面的测地线中理解。换一种解释，并不改变这句话的量词；改变的是点、线以及“位于其上”的含义。要用这种方法判断平行公设是否来自其他假设，必须先把那些假设写清楚。否则，一个图形即使没有唯一平行线，也可能同时破坏了原本希望保留的条件。

\subsection{关联、顺序、合同及连续性}
\label{g1:subsec:A01:01}

先只考虑点和线的相互归属，不讨论长度、角度或介于关系。这允许我们把模型核验压缩到一组明确的条件，并看见后来加入的结构分别做了什么。

\begin{definition}[点线关联结构]
	\label{g1:def:incidence-model}
	设 \(X\) 是点集，\(\mathcal L\) 是线集，\(I\subset X\times\mathcal L\) 是一个关系。称 \((X,\mathcal L,I)\) 为一个\term[guanlianjiegou]{关联结构}[incidence structure]；\(I(A,\ell)\) 表示点 \(A\) 位于线 \(\ell\) 上。三个点若同在某条线上，就称它们共线。本节以 \(\mathsf I\) 表示以下三个公理组成的系统：
	\symidx{\(\mathsf I\)：附录 A 的三个关联公理}
	\begin{enumerate}
		\item 任意两个不同点恰好同在一条线上。
		\item 每条线上至少有两个不同点。
		\item 存在三个不共线的点。
	\end{enumerate}
	为点、线及关联关系指定具体集合与关系，并使这三个条件全部成立，就得到 \(\mathsf I\) 的一个\term[moxing]{模型}[model]。
\end{definition}

第一条中的“任意”“存在”和“唯一”不能在核验中互相替代。一个图中画出的两点可能有连线，仍不足以说明所有两点都有连线；若点集或线集受到限制，熟悉的直线也可能根本不属于所选模型。第二、三条排除只有空线、单点线或所有点都挤在同一条线上的退化情况。

\begin{proposition}[两线的公共点]
	\label{g1:prop:incidence-intersection}
	在满足 \(\mathsf I\) 的关联结构中，两条不同的线至多有一个公共点；对每条线，至少有一个点不在该线上。
\end{proposition}
\begin{proof}
	若两条线有不同的公共点 \(A,B\)，则它们都经过 \(A,B\)，由第一条的唯一性只能是同一条线。再固定一条线 \(\ell\)，取第三条保证的三个不共线点。它们不能全在 \(\ell\) 上，故其中至少一个点不在 \(\ell\) 上。
\end{proof}

这些公理尚未规定一条线上的点怎样排列。\term[jieyuguanxi]{介于关系}[betweenness relation]用三点关系 \(B(A,B,C)\) 表示“\(B\) 严格介于 \(A,C\) 之间”；严格意味着三点互不相同。\term[xianduanhetong]{线段合同}[segment congruence]则用四点关系 \(AB\equiv CD\) 表示两段可以作为同一种长度来比较。在公理化处理时，可以把这些关系也作为原始关系，另列它们应满足的条件。关联只说明哪些点共线，顺序说明谁在谁之间，合同说明哪些段相等，三者并不是同一信息。

\ProseParagraphBreak

本卷正文已经使用距离，因此可以在模型中直接解释这些新关系。设点集 \(X\) 另带距离 \(d\)，约定
\[
	\begin{aligned}
		B(A,B,C)\iff{}& A,B,C\text{ 互异且共线，}\\
		&\text{且 }d(A,C)=d(A,B)+d(B,C),\\
		AB\equiv CD\iff{}& d(A,B)=d(C,D).
	\end{aligned}
\]
共线条件仍须保留；在一般距离空间里，三角不等式取等号并不能单独指定几何直线。下面的性质说明，把每条线参数化为通常的实数轴，能够保证哪些顺序、合同与连续性条件。

\begin{proposition}[实参数给出的线内结构]
	\label{g1:prop:axiom-line-parameter}
	设关联结构 \((X,\mathcal L,I)\) 满足 \(\mathsf I\)，点集带距离 \(d\)。假设对每条线 \(\ell\)，都有一个从 \(\mathbb R\) 到 \(\ell\) 上点集的双射 \(\gamma_\ell\)，满足
	\[
		d\bigl(\gamma_\ell(s),\gamma_\ell(t)\bigr)=|s-t|.
	\]
	按上述方式定义介于与线段合同，则有以下性质：
	\begin{enumerate}
		\item \(B(A,B,C)\) 与 \(B(C,B,A)\) 等价；三个互异共线点恰有一个介于另外两个之间。
		\item 对不同点 \(A,B\)，存在 \(C\) 使 \(B(A,B,C)\) 成立。
		\item 固定起点 \(A\) 和一条从 \(A\) 出发的射线，对任意 \(r>0\)，射线上恰有一个点 \(D\) 满足 \(d(A,D)=r\)。特别地，任意非退化线段都可以在该射线上唯一复制。
		\item 线段合同是点对上的等价关系；若 \(B(A,B,C)\)、\(B(A',B',C')\)，且 \(AB\equiv A'B'\)、\(BC\equiv B'C'\)，则 \(AC\equiv A'C'\)。
		\item 每条线选定一个方向后，其中每个非空、有上界的点集都有最小上界。
	\end{enumerate}
\end{proposition}
\begin{proof}
	固定一条线并用 \(\gamma_\ell\) 的实参数表示其点。对三个互异参数 \(a,b,c\)，等式
	\[
		|a-c|=|a-b|+|b-c|
	\]
	成立，当且仅当 \(a<b<c\) 或 \(c<b<a\)。这给出第一项；取 \(c=2b-a\) 则给出第二项。

	射线可以表示为 \(\gamma_\ell([a,\infty))\) 或 \(\gamma_\ell(({-\infty},a])\)。相应的唯一参数为 \(a+r\) 或 \(a-r\)，从而得到第三项。合同的自反、对称与传递性来自实数相等；第四项最后的断言来自两组距离的相加。

	最后按所选方向令参数递增；若原参数方向相反，改用 \(t\mapsto\gamma_\ell(-t)\)。给定非空、有上界的点集 \(S\)，其参数集 \(\gamma_\ell^{-1}(S)\) 是非空、有上界的实数集，故有最小上界 \(s_0\)。点 \(\gamma_\ell(s_0)\) 是 \(S\) 的一个上界，而任何其他上界的参数都不小于 \(s_0\)，所以它就是所需的最小上界。
\end{proof}

这里的射线、方向和顺序没有依赖某个特选参数。事实上，若 \(\gamma,\widetilde\gamma\) 是同一条线的两种上述参数，令 \(h=\gamma^{-1}\circ\widetilde\gamma\)，则 \(|h(s)-h(t)|=|s-t|\)。记 \(b=h(0)\)，由到 \(0\) 的距离知 \(h(t)-b\) 的绝对值为 \(|t|\)，再与 \(h(1)-b=1\) 或 \(-1\) 比较到 \(1\) 的距离，分别得到 \(h(t)=b+t\) 或 \(h(t)=b-t\)。故换参数至多改变原点和方向，不改变介于关系，也不改变一个方向上最小上界是否存在。

\ProseParagraphBreak

第五项是一种\term[xuwanbeixing]{序完备性}[order completeness]要求。它与前四项的区别在于，需要控制无限点集留下的空隙；证明实际调用了实数的最小上界性质。若只允许变量遍历点和线，“任意点集”就不能直接写成这两类变量的一次全称量化。要把这条要求纳入形式公理系统，必须另行说明采用集合变量，还是采用某种较弱的一阶公理方案。下文讨论一阶独立性时，所用公理系统仍是明确列出的 \(\mathsf I\)。

\ProseParagraphBreak

也不能把上面的五项性质当成整个平面几何的公理表。它们没有规定一条线穿过三角形一边之后如何与其他边相交，也没有给出角合同以及边角边条件。Hilbert 的《几何基础》把这些不同任务分组处理；本附录所据 Townsend 英译重印本的关联、顺序、平行、合同、连续性依次见\cite[Chapter I, \S\S 2--8]{Hilbert1950Foundations}。其中 \S 8 以等段反复截取的 Archimedes 条件开头，随后另论完备性，不能把不同连续性要求视为同一个句子。这里选择较小的系统，是为了让每一个模型判断都有可以逐项检查的假设。

\subsection{平行公设和非欧模型}
\label{g1:subsec:A01:02}

在关联结构中，可以只用“有没有公共点”定义不相交，而无须先定义角度。对点线组成的整个结构，考虑如下\term[pingxinggongse]{平行公设}[parallel postulate]：
\[
	\begin{gathered}
		\mathsf P:\quad\text{对每条线 \(\ell\) 和每个不在 \(\ell\) 上的点 \(A\)，}\\
		\text{恰有一条过 \(A\) 的线与 \(\ell\) 不相交。}
	\end{gathered}
\]
这里采用的是唯一平行线的表述，与上述英译本\cite[Chapter I, \S 5]{Hilbert1950Foundations}一致。“不相交”指在模型的点集中没有公共点。命题 \(\mathsf P\) 同时包含存在性和唯一性；否定它，可以是某组点线根本没有这样的线，也可以是至少有两条。

\begin{proposition}[欧氏平面的模型核验]
	\label{g1:prop:euclidean-incidence-model}
	取点集 \(X=\mathbb R^2\)，线集为全部仿射直线 \(a+\mathbb R u\)，其中 \(u\ne0\)，关联解释为点属于直线的关系。这个结构满足 \(\mathsf I\) 和 \(\mathsf P\)。配以欧氏距离后，每条线满足\cref{g1:prop:axiom-line-parameter}的参数条件。
\end{proposition}
\begin{proof}
	两不同点 \(a,b\) 确定直线 \(a+\mathbb R(b-a)\)。若另一条直线包含它们，其方向空间必包含非零向量 \(b-a\)，因方向空间是一维的，它只能等于 \(\mathbb R(b-a)\)，故两线相同。每条线包含 \(a\) 和 \(a+u\)；点 \((0,0),(1,0),(0,1)\) 不共线。因此三条关联公理成立。

	固定 \(\ell=a+\mathbb R u\) 和 \(p\notin\ell\)。线 \(m=p+\mathbb R u\) 与 \(\ell\) 不交，否则 \(p-a\in\mathbb R u\)，与 \(p\notin\ell\) 矛盾。若过 \(p\) 的另一条线方向为 \(v\)，且 \(v\) 与 \(u\) 线性无关，则 \(u,v\) 是 \(\mathbb R^2\) 的一组基，所以方程 \(a+su=p+tv\) 有解，两线相交。因而不相交的线只能取方向 \(\mathbb R u\)，也就只能是 \(m\)。最后将 \(u\) 归一化为单位向量，\(\gamma(t)=a+tu\) 就是所需参数。
\end{proof}

同样的关联公理也允许另一种解释。取上半平面
\[
	\mathbb H=\{x+\mathrm i y\in\mathbb C:y>0\},
\]
把垂直于实轴的直线，以及圆心在实轴上的开上半圆，全部称为模型的“线”。\cref{g1:thm:hyperbolic-distance-geodesics}已证明，这些正是度量 \((\mathrm dx^2+\mathrm dy^2)/y^2\) 的完整测地线，每对点间的测地线段唯一且实现距离。现在先核验点线关联；这一步甚至还不需要距离公式。

\begin{proposition}[双曲平面的关联与平行]
	\label{g1:prop:hyperbolic-incidence-model}
	以上半平面 \(\mathbb H\) 为点集，以竖直线和实轴为圆心的开上半圆为线集，并以点属于曲线解释关联，得到 \(\mathsf I\) 的一个模型。该模型不满足 \(\mathsf P\)。
\end{proposition}
\begin{proof}
	取不同点 \(z_j=x_j+\mathrm i y_j\)，其中 \(y_j>0\)。若 \(x_1=x_2\)，竖直线 \(x=x_1\) 经过两点；一个圆心为 \((c,0)\) 的圆若也经过两点，则
	\[
		(x_1-c)^2+y_1^2=(x_2-c)^2+y_2^2
	\]
	迫使 \(y_1=y_2\)，与两点不同矛盾。因此竖直线是唯一选择。若 \(x_1\ne x_2\)，竖直线不可能同时经过两点，上式唯一确定圆心
	\[
		c=\frac{x_1^2+y_1^2-x_2^2-y_2^2}{2(x_1-x_2)}.
	\]
	半径 \(r=\sqrt{(x_1-c)^2+y_1^2}>0\) 随之确定，并且所得上半圆确实经过两点。这证明了第一条关联公理。每条竖直线和每个开上半圆都包含无穷多个点；而 \(\mathrm i,2\mathrm i,1+\mathrm i\) 不共线，因为前两点确定竖直线 \(x=0\)，第三点不在其上。

	为否定 \(\mathsf P\)，只需找到一组点线。令
	\[
		\ell=\{\mathrm i y:y>0\},\qquad p=1+\mathrm i.
	\]
	过 \(p\) 的竖直线 \(m_1=\{1+\mathrm i y:y>0\}\) 与 \(\ell\) 不交。开上半圆
	\[
		m_2=\{x+\mathrm i y:(x-1)^2+y^2=1,\ y>0\}
	\]
	也过 \(p\)，而且同样与 \(\ell\) 不交：令 \(x=0\) 后圆方程只允许 \(y=0\)，这不属于 \(\mathbb H\)。由于 \(m_1\ne m_2\)，唯一性失败。
\end{proof}

\begin{figure}[htbp]
	\centering
	% !TeX root = ../Book1.tex
\begin{tikzpicture}[>=stealth,font=\small,x=1.35cm,y=1.35cm]
	\begin{scope}
		\draw[thick,<->] (0,-.3)--(0,2.25) node[above] {\(\ell\)};
		\draw[blue!70!black,thick,<->] (1,-.3)--(1,2.25) node[above] {\(m\)};
		\fill (1,1) circle (1.5pt) node[right] {\(p\)};
		\node at (.5,-.65) {欧氏平面};
	\end{scope}
	\begin{scope}[xshift=4.5cm]
		\draw[dashed,gray,->] (-.35,0)--(2.5,0) node[right] {\(\mathbb R\)};
		\draw[thick,->] (0,0)--(0,2.25) node[above] {\(\ell\)};
		\draw[blue!70!black,thick,->] (1,0)--(1,2.25) node[above] {\(m_1\)};
		\draw[red!70!black,thick] (0,0) arc[start angle=180,end angle=0,radius=1];
		\node[red!70!black,right] at (1.8,.65) {\(m_2\)};
		\fill (1,1) circle (1.5pt) node[above right] {\(p=(1,1)\)};
		\foreach \x in {0,1,2} {
			\draw[gray,fill=white] (\x,0) circle (1.6pt);
		}
		\node[below] at (0,0) {\(0\)};
		\node[below] at (2,0) {\(2\)};
		\node at (1,-.65) {上半平面模型};
	\end{scope}
\end{tikzpicture}

	\caption[平行关系依赖所选模型]{左图在欧氏平面中，过 \(p\) 与 \(\ell\) 不交的直线只有 \(m\)。右图在上半平面模型中，\(m_1\) 与 \(m_2\) 是两条不同的模型直线，都经过 \(p\) 且不与 \(\ell\) 相交；空心端点位于被排除的实轴上。图只显示线的一部分。}
	\label{g1:fig:axiom-parallel-models}
\end{figure}

\cref{g1:fig:axiom-parallel-models}中，右侧半圆在图纸上弯曲，在模型中却是一条直线。把图纸上的欧氏直线含义强加给它，会混淆解释模型所用的环境与模型自身的关系。同样，\((0,0)\) 只是 \(\ell\) 和 \(m_2\) 的共同理想端点；它不是模型中的点，因而不能充当交点。

\begin{proposition}[双曲直线的实参数]
	\label{g1:prop:hyperbolic-line-parameter}
	以上半平面的双曲距离度量点集，并以竖直线及实轴为圆心的开上半圆为线，则每条线都有\cref{g1:prop:axiom-line-parameter}要求的等距参数 \(\mathbb R\rightarrow\ell\)。因此该模型具有该命题列出的介于、线段复制、合同相加及序完备性性质。
\end{proposition}
\begin{proof}
	对竖直线 \(x=c\)，取 \(\gamma(t)=c+\mathrm i e^t\)。它是从 \(\mathbb R\) 到该线的双射。由\eqref{g1:eq:half-plane-distance}，
	\[
		\cosh d\bigl(\gamma(s),\gamma(t)\bigr)
		=1+\frac{(e^s-e^t)^2}{2e^{s+t}}
		=\cosh(s-t).
	\]
	由于距离非负，且 \(\cosh\) 在 \([0,\infty)\) 上单调递增，距离就是 \(|s-t|\)。

	再取理想端点为实数 \(a<b\) 的上半圆。变换
	\[
		F(z)=\frac{z-a}{b-z}
	\]
	对应的实矩阵为 \(\begin{psmallmatrix}1&-a\\-1&b\end{psmallmatrix}\)，其行列式为 \(b-a>0\)；乘以 \((b-a)^{-1/2}\) 后行列式为 \(1\)，而分式不变。由\cref{g1:thm:hyperbolic-isometry-group}中已直接验证的度量保持公式，\(F\) 是上半平面的双曲等距。具体地，
	\[
		\operatorname{Im}F(z)=\frac{(b-a)y}{|b-z|^2},\qquad
		|F'(z)|=\frac{b-a}{|b-z|^2},
	\]
	所以长度元 \(|\mathrm dz|/y\) 保持不变；逆变换是 \(F^{-1}(w)=(bw+a)/(1+w)\)。圆方程为 \(y^2=(x-a)(b-x)\)，而
	\[
		\operatorname{Re}F(z)=\frac{(x-a)(b-x)-y^2}{|b-z|^2}.
	\]
	因此 \(F\) 将该开上半圆双射到正虚轴。用 \(\gamma(t)=F^{-1}(\mathrm i e^t)\) 即得所需参数，距离保持性将竖直线的结论传回原来的半圆。
\end{proof}

这个核验说明，双曲模型中并非连线段延长、长度复制和线内连续性也一并失效。两种模型在线内满足上述共同性质，却有不同的平行关系。但是，这仍只证明了已经列出的性质；如果把公理系统加强到包含所有平面顺序与角合同公理，就需要相应增加模型核验。

\subsection{独立性与相对一致性}
\label{g1:subsec:A01:03}

现在把模型比较转成关于证明的结论。设 \(T\) 是一组公理，\(\varphi\) 是用同一种语言写出的命题。记
\[
	\mathcal M\models T
\]
表示结构 \(\mathcal M\) 中 \(T\) 的每条公理都成立。例如，在本节的语言中，“任意点”遍历 \(\mathcal M\) 的点集，“存在一条线”遍历它的线集，原始关系 \(I\) 由模型指定的关联关系解释。这称为\term[manzuguanxi]{满足关系}[satisfaction relation]。\symidx{\(\mathcal M\models T\)：模型满足公理系统}记 \(T\vdash\varphi\) 表示从 \(T\) 出发，按照固定的一阶逻辑推演规则，经过有限步得到 \(\varphi\)；它说的是存在一份形式证明。\symidx{\(T\vdash\varphi\)：从公理系统形式推导命题}

\ProseParagraphBreak

两种记号之间，当前需要的是“证明不会把模型中的真命题变成假命题”这一方向。为看清它的依据，可考虑以逻辑公理、等式规则、分离规则和全称化组成的通常一阶推演：分离规则由 \(\alpha\) 及 \(\alpha\Rightarrow\beta\) 推出 \(\beta\)；全称化由对任意变量赋值均成立的 \(\alpha(x)\) 推出 \(\forall x\,\alpha(x)\)，且该变量不能在尚未解除的额外假设中自由出现。等式规则只允许用相等对象替换同一表达式中的对象。逻辑公理包括命题逻辑恒真式以及合法代入和量词分配等公理，它们的真值由逻辑连接词、量词与相等的解释确定，与选哪一个几何模型无关。

\begin{proposition}[推导的保真性]
	\label{g1:prop:geometric-soundness}
	设 \(T\) 是由一阶句子组成的公理系统，推导使用上述通常的一阶逻辑规则。若 \(\mathcal M\models T\) 且 \(T\vdash\varphi\)，则 \(\mathcal M\models\varphi\)。
\end{proposition}
\begin{proof}
	按一份有限证明中各行的次序归纳，证明每行在 \(\mathcal M\) 的任意变量赋值下成立。若该行是 \(T\) 的公理，它由模型假设成立；若是逻辑公理，它由逻辑语义成立。例如 \(\forall x\,\alpha(x)\Rightarrow\alpha(t)\) 在代入不使变量被错误绑定时成立，是因为全称命题包含了由项 \(t\) 指定的那个对象。量词分配
	\[
		\forall x\,(\alpha\Rightarrow\beta)
		\Rightarrow\bigl(\alpha\Rightarrow\forall x\,\beta\bigr)
	\]
	在 \(x\) 不自由出现在 \(\alpha\) 中时成立，是因为改变 \(x\) 的赋值不改变 \(\alpha\) 的真值。等式的自反性与替换规则来自相同对象具有相同解释。

	若该行由分离规则得到，归纳假设已知 \(\alpha\) 和 \(\alpha\Rightarrow\beta\) 对同一赋值为真，故 \(\beta\) 为真。若由全称化得到，归纳假设说明前一公式对 \(x\) 的每个允许取值都为真，因而其全称量化为真。这里 \(T\) 由句子组成，没有自由变量妨碍全称化；若在带临时假设的证明中使用该规则，就必须保留前述自由变量限制。由归纳，最后一行 \(\varphi\) 在模型中成立。
\end{proof}

\begin{definition}[独立性与一致性]
	\label{g1:def:axiom-independence}
	给定公理系统 \(T\) 与固定的推演规则，若 \(T\nvdash\varphi\) 且 \(T\nvdash\neg\varphi\)，则称命题 \(\varphi\) 对 \(T\) 是\term[dulixing]{独立的}[independent]。若不存在一个句子 \(\psi\) 使 \(T\vdash\psi\) 与 \(T\vdash\neg\psi\) 同时成立，则称 \(T\) 是\term[yizhixing]{一致的}[consistent]。
\end{definition}

\begin{theorem}[本节关联公理中的平行独立性]
	\label{g1:thm:incidence-parallel-independence}
	对\cref{g1:def:incidence-model}的三个关联公理 \(\mathsf I\)，唯一平行线命题 \(\mathsf P\) 是独立的。系统 \(\mathsf I\cup\{\mathsf P\}\) 和 \(\mathsf I\cup\{\neg\mathsf P\}\) 都有本节构造的模型。
\end{theorem}
\begin{proof}
	\cref{g1:prop:euclidean-incidence-model}给出一个满足 \(\mathsf I\) 与 \(\mathsf P\) 的模型，而\cref{g1:prop:hyperbolic-incidence-model}给出一个满足 \(\mathsf I\) 与 \(\neg\mathsf P\) 的模型。若 \(\mathsf I\vdash\mathsf P\)，\cref{g1:prop:geometric-soundness}便迫使双曲模型满足 \(\mathsf P\)，与已核验的两条不交线矛盾。若 \(\mathsf I\vdash\neg\mathsf P\)，同一保真性又迫使欧氏模型满足 \(\neg\mathsf P\)，与欧氏模型中存在且唯一的平行线矛盾。这排除了两个方向的形式证明。
\end{proof}

同样的论证解释了模型存在为何能支持一致性。若某个有模型的公理系统同时推出 \(\psi\) 与 \(\neg\psi\)，保真性便使同一模型中的同一句子既真又假，这是不可能的。这里没有使用“一切在所有模型中成立的命题都可证明”的反方向，也就没有调用逻辑完备性定理。

\ProseParagraphBreak

不过，我们用来构造模型的实数、集合、函数及有限推导本身，也处在一个背景理论中。上述证明说明：只要背景理论中的这些构造和推理有效，\(\mathsf I\) 的两个扩充就各有模型，因而不可能产生形式矛盾。这就是当前的\term[xiangduiyizhixing]{相对一致性}[relative consistency]结论所依赖的背景。它没有脱离一切数学前提证明一个绝对的“永不矛盾”；它把待研究的公理系统与已经承担构造任务的背景数学联系起来。

\ProseParagraphBreak

最后，独立性总是相对于写明的公理系统而言。从 \(\mathsf I\) 推不出 \(\mathsf P\)，并不自动说明向 \(\mathsf I\) 加入任意新公理后仍推不出它。经典几何基础中更强的任务，是保留选定体系的其他全部公理，再说明平行公设不能推出；这要求非欧模型也满足那些顺序、合同及连续性条件。Hilbert 在所据版本\cite[Chapter II, \S\S 9--10]{Hilbert1950Foundations}中先讨论公理的相容性，再讨论平行公理的独立性。阅读这种论证时，应把所保留的公理表与模型解释逐项对照；本节三个关联公理的定理则提供了一次可以直接检查的较小规模实例。

\section{坐标化及历史论证}
\label{g1:sec:A02}

把几何对象写成数对，可以将构造问题变成方程；但数对与原对象之间的对应仍须建立。先从已经给定的实仿射空间检验这项对应，再回到一项古典作图，就能分清两件事：现代坐标证明是否成立，以及它与原始证明实际保留了多少相同的假设和步骤。

\subsection{几何对象与代数模型的双向建立}
\label{g1:subsec:A02:01}

给平面上的点写上两个实数，并不会自动证明一个公理系统中的每个平面都能这样表示。坐标化至少包含两个不同的问题：已经有向量空间和平移运算时，怎样用坐标忠实地记录几何关系；只给抽象的点、线与若干关系时，又能否从这些关系中构造数域及向量运算。\chapref{g1:ch:01}解决的是前一个问题，其出发点在\cref{g1:def:affine-space}中已经写明：方向空间是有限维实向量空间，且任意两点之间有唯一的平移向量。本小节把这个来回构造及其条件展开，以便准确判断“化为代数问题”究竟保留了什么。

\ProseParagraphBreak

设 \(A\) 为 \(n\) 维实仿射空间，方向空间为 \(V\)，其中 \(n\ge1\)。选一个原点 \(o\in A\) 和 \(V\) 的基 \(e_1,\ldots,e_n\)，记这组选择为 \(\mathcal F\)。\cref{g1:prop:affine-coordinates}中的标架是 \((o,o+e_1,\ldots,o+e_n)\)。该命题的存在与唯一性给出互逆映射
\[
	\begin{aligned}
		\Phi_{\mathcal F}:A&\longrightarrow\mathbb R^n,
		& p=o+\sum_i x_ie_i&\longmapsto x,\\
		\Psi_{\mathcal F}:\mathbb R^n&\longrightarrow A,
		& x&\longmapsto o+\sum_i x_ie_i.
	\end{aligned}
\]
这是点集之间的双射。要说它们是几何结构的同构，还必须检查直线、介于关系和所要保留的其他关系怎样对应。

\begin{proposition}[坐标中的几何关系]\label{g1:prop:affine-model-relations}
	设 \(A\) 为 \(n\) 维实仿射空间，\(\mathcal F=(o;e_1,\ldots,e_n)\) 为一组仿射标架数据。若 \(p,q,r\) 的坐标分别为 \(x,y,z\)，则有以下对应。
	\begin{enumerate}
		\item 对 \(p\ne q\)，过 \(p,q\) 的直线在坐标中恰为 \(x+\mathbb R(y-x)\)。因此点在线上的关联关系被 \(\Phi_{\mathcal F}\) 及其逆保持。
		\item 对 \(p\ne q\)，\(r\) 严格介于 \(p,q\) 之间，当且仅当存在 \(0<t<1\) 使 \(z=(1-t)x+ty\)。这里严格介于的定义是 \(p\ne q\) 且 \(r=p+t(q-p)\)、\(0<t<1\)。
		\item 若 \(V\) 另有给定的正定内积，记 \(G=(\langle e_i,e_j\rangle)\)，则 \(G\) 对称正定，且
		\begin{equation}\label{g1:eq:coordinate-model-distance}
			d(p,q)^2=(y-x)^{\mathsf T}G(y-x).
		\end{equation}
		所以线段 \(pq\) 与 \(p'q'\) 合同，等价于它们的坐标差对二次型 \(G\) 具有相同取值。对以 \(p\) 为顶点、方向非零的两条射线 \(pq,pr\)，夹角 \(\theta\in[0,\pi]\) 满足
		\[
			\cos\theta=
			\frac{(y-x)^{\mathsf T}G(z-x)}
			{\sqrt{(y-x)^{\mathsf T}G(y-x)}\,
			\sqrt{(z-x)^{\mathsf T}G(z-x)}}.
		\]
	\end{enumerate}
\end{proposition}
\begin{proof}
	令 \(E:\mathbb R^n\to V\) 为 \(Ex=\sum_i x_ie_i\)。它是线性同构，且点差满足 \(q-p=E(y-x)\)。因而
	\[
		p+t(q-p)=o+E\bigl(x+t(y-x)\bigr).
	\]
	令 \(t\) 遍历全部实数，得到直线的精确对应；再将参数限制为 \(0<t<1\)，得到严格介于关系的对应。由于坐标映射双射，这些关系在两个方向都被保持。

	若有给定内积，双线性使 \(\langle Eu,Ev\rangle=u^{\mathsf T}Gv\)。矩阵 \(G\) 的对称性来自内积对称性；对非零 \(u\)，\(Eu\ne0\)，所以 \(u^{\mathsf T}Gu=\langle Eu,Eu\rangle>0\)。把 \(u=y-x\) 代入就得到距离公式。线段合同指两段距离相等，而距离非负，故等价于平方相等。夹角公式则是\cref{g1:def:inner-product}之后给出的内积夹角定义在此坐标中的表达；两射线方向非零使分母严格为正。\(\cos\) 在 \([0,\pi]\) 上单射，所以比较夹角也可等价地比较右侧数值。
\end{proof}

反方向从代数数据开始同样具体。给定 \(\mathbb R^n\) 和对称正定实矩阵 \(G\)，以坐标加法作平移，令直线为 \(x+\mathbb Rv\)、\(v\ne0\)，令严格介于由 \(0<t<1\) 的仿射组合定义，并令 \(\langle u,v\rangle_G=u^{\mathsf T}Gv\)。这满足\cref{g1:def:affine-space,g1:def:inner-product}的条件，因而给出一个实欧氏仿射空间。对原来的 \(A\) 按此办法重建后，\(\Psi_{\mathcal F}\) 恢复原点集、点差、直线、介于及内积，而不只给出一幅外形相似的图。若没有指定 \(G\)，相同的坐标数据只给出仿射结构及其基于实数次序的介于关系。

\ProseParagraphBreak

各项假设在这里承担不同任务。实向量空间给出加法与实数倍，所以能用一组基坐标化；实数的次序使 \(0<t<1\) 有意义；正定内积才给出欧氏长度和夹角。仅就仿射组合及介于的公式而言，可以把 \(\mathbb R\) 换成一个已经给定的有序域，但相应平方根是否存在、直线是否具备实数的连续性，还需另外讨论。不能因为写出了同样的式子，就把这些额外性质一起带入。

\ProseParagraphBreak

换一组标架会怎样？设新原点的旧坐标为 \(a\)，新基在旧基中的列矩阵为 \(S\in\operatorname{GL}(n,\mathbb R)\)。\cref{g1:prop:affine-coordinates}给出同一点的新旧坐标 \(x=a+Sx'\)，故
\begin{equation}\label{g1:eq:coordinate-model-change}
	x'=S^{-1}(x-a),\qquad G'=S^{\mathsf T}GS.
\end{equation}
第二式由新基向量的内积直接得到。下面的交换图表示先取旧坐标再换坐标，与直接取新坐标相同；左边始终是同一个几何点。
\[
	\begin{tikzcd}[column sep=large,row sep=large]
		A \arrow[r,"{\Phi_{\mathcal F}}"] \arrow[d,equal]
		& (\mathbb R^n,G) \arrow[d,"{x\mapsto S^{-1}(x-a)}"]\\
		A \arrow[r,"{\Phi_{\mathcal F'}}"']
		& (\mathbb R^n,G')
	\end{tikzcd}
\]
对点差 \(u=Su'\)，有 \(u^{\mathsf T}Gu=(u')^{\mathsf T}G'u'\)，所以两套坐标计算同一个距离。\(G\) 与 \(G'\) 的数值可以不同，几何长度并没有改变。特别地，非正交标架中的坐标差一般不能用各分量的平方和计算长度。

\ProseParagraphBreak

同一计算也解释了变换群怎样随保留结构而改变。在固定坐标模型中，\(x\mapsto Mx+b\)、\(M\in\operatorname{GL}(n,\mathbb R)\)，保持直线、介于及全部仿射组合；它保持 \(G\) 给出的数值距离，当且仅当 \(M^{\mathsf T}GM=G\)。必要性可先比较每个向量的长度平方，再用\cref{g1:prop:polarization}恢复两个向量的配对；充分性则由距离公式直接得到。标准正交坐标中的条件就是 \(M^{\mathsf T}M=I\)。这里只在描述一组已经证明为仿射的映射；若一开始只给一个集合映射，就不能未经论证把它写成矩阵加平移。

\begin{proposition}[连续的中点保持映射]\label{g1:prop:continuous-midpoint-affine}
	设 \(A,B\) 为有限维实仿射空间，赋予任一仿射坐标确定的通常拓扑。若映射 \(f:A\to B\) 在至少一个点连续，且对任意 \(p,q\in A\) 满足
	\[
		f\bigl(\tfrac12p+\tfrac12q\bigr)
		=\tfrac12 f(p)+\tfrac12 f(q),
	\]
	则 \(f\) 为实仿射映射。反之，每个实仿射映射都连续且保持中点。
\end{proposition}
\begin{proof}
	取 \(A\) 的一个原点 \(o\)，以 \(f(o)\) 为 \(B\) 的原点，并选两边方向基。于是只需考察一个保持中点、满足 \(F(0)=0\) 的映射 \(F:\mathbb R^n\to\mathbb R^m\)。中点条件对 \(0,2x\) 给出 \(F(2x)=2F(x)\)，所以
	\[
		F(x+y)=2F\bigl((x+y)/2\bigr)=F(x)+F(y).
	\]
	由加法性可依次得到整数倍以及有理数倍相容性，即 \(F(rx)=rF(x)\)、\(r\in\mathbb Q\)。若连续点为 \(x_0\)，则当 \(h\to0\) 时，\(F(h)=F(x_0+h)-F(x_0)\to0\)，所以 \(F\) 在零点连续；再由 \(F(x+h)-F(x)=F(h)\)，它在每一点都连续。

	对任意 \(t\in\mathbb R\)，取有理数列 \(r_j\to t\)。连续性使
	\[
		F(tx)=\lim_jF(r_jx)=\lim_jr_jF(x)=tF(x).
	\]
	因此 \(F\) 实线性，恢复原点便得 \(f\) 仿射。反之，仿射映射在坐标中为有限维线性映射加常向量，因而连续，并由\cref{g1:prop:affine-preservation}保持所有仿射组合，特别保持中点。
\end{proof}

这个命题把\cref{g1:thm:isometry-affine}证明中用到的恢复步骤分离出来，明确了连续性的作用：中点等式先恢复有理数倍，实数倍还须由极限恢复。对欧氏空间的距离保持映射，连续性来自距离等式，中点保持来自距离对中点的唯一刻画，所以那一定理能够进一步推出仿射性；它没有预先把仿射性放进“等距”的定义。

\begin{example}[关联和介于不足以确定合同]\label{g1:exm:incidence-order-not-congruence}
	在同一个点集 \(\mathbb R^2\) 上，使用同一族通常直线及同一介于关系，分别取
	\[
		G_1=\begin{pmatrix}1&0\\0&1\end{pmatrix},
		\qquad G_2=\begin{pmatrix}4&0\\0&1\end{pmatrix}.
	\]
	两种结构的点线关联及介于关系完全相同，但线段 \((0,0)(1,0)\) 与 \((0,0)(0,1)\) 在第一种结构中合同，在第二种结构中长度分别为 \(2,1\)，并不合同。因此关联和介于数据没有唯一指定合同关系。两模型作为未标记的欧氏空间仍然等距：\((x_1,x_2)\mapsto(2x_1,x_2)\) 将第二种距离变成第一种距离。这里不同的是同一批指定点之间保留了哪些关系，不能把这一差别误说成存在两种非同构的二维实欧氏空间。

	若只把 \(G_1\) 改为 \(cG_1\)、\(c>0\)，则所有长度同时乘 \(\sqrt c\)，线段之间是否合同和夹角大小都不变。因此即使给出了全部合同关系，数值长度仍需要选定单位。把“保持合同关系”与“逐个保持给定的数值距离”混为一谈，也会错误地删去统一伸缩这一自由度。
\end{example}

\chapref{g1:ch:02}的射影模型有同样需要说明的出发点。\cref{g1:def:projective-space}先给定实或复向量空间 \(V\)，再把一维子空间当作点；\cref{g1:prop:projective-charts}在一个非零齐次坐标处归一化，得到仿射图。这建立了向量空间与所构造射影空间之间的对应，却没有从任意抽象点线关联体系中推出底层数域。对于后一问题，域运算怎样定义、是否满足域公理、几何关系是否全部恢复，都需要专门的条件与证明。分清这两种坐标化任务，才能知道代数模型是在实现给定公理、描述一个已经构造的空间，还是承担某个几何对象的分类。

\subsection{历史证明的原始语境与现代重述}
\label{g1:subsec:A02:02}

一幅图中两条圆周相交，与一套公理能够保证它们相交，是两个需要分别说明的事实。阅读历史证明时，还须追问所读的文字属于哪个版本：是原著的译文，后人的编订，还是加入新公理后的重写。下面以一个等边三角形作图和两种平行条目为例，把文本中实际写出的论证，与本书在实坐标中补出的论证并排考察。

\ProseParagraphBreak

这里通过后世编订本考察 Euclid 的《几何原本》。\footnote{欧几里得（\OriginalName{Εὐκλείδης}，拉丁转写 Eukleidēs，约公元前300年活动，确切生卒年不详），古希腊数学家，活动于亚历山大里亚，以《几何原本》对几何与数论的系统论述著称。}具体使用的是 Casey 于1885年出版的第三版。\footnote{John Casey（约翰·凯西，1820--1891），爱尔兰数学家，研究初等几何与圆的几何，编写过《Euclid 几何原本》及其续编。}它在题名中明确列出注释与习题，并保留了大量编者说明，因此以下只把它当作十九世纪的一个可核查文本，不把它逐字归给古希腊原著。该版 Book I 的 Prop. I 要求在给定的有限直线上构造等边三角形；其中的 \emph{finite right line} 指两端已给定的直线段，\emph{right line} 在这里是直线，并非带有直角条件的线。该版在 Postulates, I 后也明确解释了这个用语。\cite[Book I, Postulates, I; Prop. I]{Casey1885Euclid}

\ProseParagraphBreak

把已给线段记为 \(AB\)，该文本的构造分为三个动作：以 \(A\) 为圆心、\(AB\) 为半径作圆；以 \(B\) 为圆心、\(BA\) 为半径作第二个圆，并取交点 \(C\)；连接 \(CA,CB\)。随后利用同圆半径相等得到 \(AC=AB\) 与 \(BC=BA\)，再用“等于同一量的量彼此相等”得到三边相等。文本对作圆和连接两点分别援引第三、第一公设，却没有在取交点处另给一个保证交点存在的论证。它证明了交点一旦取得便给出所求三角形；若要把这段论证改写成明确的现代公理演绎，交点的存在及其不在直线 \(AB\) 上，还须从所选公理推出。

\ProseParagraphBreak

在本书已经构造的实欧氏平面中，可以直接补上这一步。先选择一个定向，令 \(J\) 表示向该正向旋转 \(\pi/2\) 的线性映射。若 \(v=B-A\ne0\)，则 \(Jv\perp v\)、\(\lVert Jv\rVert=\lVert v\rVert\)。为了找同时到 \(A,B\) 等距的点，先取中点，再沿垂直方向移动。所需移动量由勾股关系确定为边长的 \(\sqrt3/2\) 倍。

\begin{proposition}[两个圆的交点与等边三角形]\label{g1:prop:equilateral-coordinate-reconstruction}
	设 \(A,B\) 是实欧氏平面中两个不同的点，固定该平面的一个定向，并令 \(J\) 为正向直角旋转。以 \(A,B\) 为圆心、以 \(\lVert B-A\rVert\) 为共同半径的两个圆周，恰好交于
	\begin{equation}\label{g1:eq:equilateral-coordinate-intersections}
		C_\pm=\frac{A+B}{2}\pm\frac{\sqrt3}{2}J(B-A).
	\end{equation}
	两个交点都不在直线 \(AB\) 上，且 \(ABC_+\)、\(ABC_-\) 都是等边三角形。
\end{proposition}
\begin{proof}
	记 \(v=B-A\)、\(r=\lVert v\rVert>0\)、\(M=(A+B)/2\)。因 \(v,Jv\) 为一组正交基，任意点 \(C\) 唯一写成 \(M+sv+tJv\)。到两个圆心的距离平方分别为
	\[
		\begin{aligned}
			\lVert C-A\rVert^2&=r^2\bigl((s+\tfrac12)^2+t^2\bigr),\\
			\lVert C-B\rVert^2&=r^2\bigl((s-\tfrac12)^2+t^2\bigr).
		\end{aligned}
	\]
	两式都等于 \(r^2\) 时，相减得 \(2s=0\)，代回得 \(t^2=3/4\)。实数中恰有 \(t=\pm\sqrt3/2\) 两解，所以交点恰为所列两点。它们的 \(Jv\) 分量非零，因而不在直线 \(AB\) 上；两个圆周的方程又给出 \(AC_\pm=BC_\pm=AB\)，即得等边性。
\end{proof}

% !TeX root = ../Book1.tex
\begin{figure}[htbp]
	\centering
	\begin{tikzpicture}[scale=2,font=\small]
		\draw[GeometryGray] (0,0) circle[radius=1];
		\draw[GeometryGray] (1,0) circle[radius=1];
		\fill[GeometryBlue!8] (0,0)--(1,0)--(0.5,{sqrt(3)/2})--cycle;
		\draw[thick,GeometryBlue] (0,0)--(1,0)--(0.5,{sqrt(3)/2})--cycle;
		\draw[GeometryGray,dashed] (0.5,{-sqrt(3)/2})--(0.5,{sqrt(3)/2});
		\draw[GeometryGray,dashed] (0,0)--(0.5,{-sqrt(3)/2})--(1,0);
		\foreach \x/\y in {0/0,1/0,0.5/0}{\fill (\x,\y) circle[radius=0.025];}
		\fill[GeometryBlue] (0.5,{sqrt(3)/2}) circle[radius=0.025];
		\fill[GeometryBlue] (0.5,{-sqrt(3)/2}) circle[radius=0.025];
		\node[below left] at (0,0) {\(A\)};
		\node[below right] at (1,0) {\(B\)};
		\node[below right] at (0.5,0) {\(M\)};
		\node[above] at (0.5,{sqrt(3)/2}) {\(C_+\)};
		\node[below] at (0.5,{-sqrt(3)/2}) {\(C_-\)};
		\node[right] at (0.5,0.44) {\(h\)};
	\end{tikzpicture}
	\caption[等边三角形作图中的两个交点]{在实欧氏平面中，两圆的共同点位于 \(AB\) 的中垂线上，距中点 \(M\) 的长度为 \(h=\sqrt3\lVert B-A\rVert/2\)。图画出的是实数模型；若改变点集，图中可见的位置未必仍是该模型中的点。}
	\label{g1:fig:equilateral-historical-construction}
\end{figure}


\cref{g1:fig:equilateral-historical-construction}把这两个交点画在实平面中。命题的证明使用了实数域、向量运算、内积和正平方根，故它是在本书所用模型中对作图的现代重述。它也说明相同的公式在适当的有序子域上仍可能有效，但并没有把这些代数与度量结构从一套更少的抽象公理中重新构造出来。若研究的问题是某组关联、顺序及合同公理是否足够，必须继续追查这些运算与交点存在性的公理依据。

\begin{example}[两个圆周分别存在而没有共同点]\label{g1:exm:rational-circles-no-intersection}
	把点集改为 \(\mathbb Q^2\)，仅考察由有理系数方程定义的两个集合
	\[
		\begin{aligned}
			S_A&=\{(x,y)\in\mathbb Q^2:x^2+y^2=1\},\\
			S_B&=\{(x,y)\in\mathbb Q^2:(x-1)^2+y^2=1\}.
		\end{aligned}
	\]
	它们都非空，例如 \((1,0)\in S_A\)、\((0,0)\in S_B\)。但若某点同时属于两者，相减得 \(x=1/2\)，随后得到 \(y^2=3/4\)，所以 \(2y\) 必须是 \(3\) 的有理平方根。这不可能：若互素整数 \(m,n\)、\(n\ne0\) 满足 \(m^2=3n^2\)，则 \(3\mid m\)，代回又得 \(3\mid n\)，矛盾。因此这两个集合不相交。

	这个计算只反驳“定义了两个圆周，就已保证它们具有所画出的交点”这一推断，并未证明某套完整几何公理中的独立性。事实上，\(\mathbb Q^2\) 配通常的距离相等关系就不满足任意射线上的线段搬移：从 \((0,0)\) 沿正 \(x\) 轴，无法取得与线段 \((0,0)(1,1)\) 等长的有理点，因为该点横坐标必须为 \(\sqrt2\)。所以不能把本例直接称作满足全部关联、顺序和合同公理的平面。
\end{example}

不同版本对平行条件的组织也不能只凭熟悉的名称判断。Casey 第三版把作图要求列为三个 Postulates，而把一个平行条件列在 Axioms, XII：若两条直线与第三条直线相遇，所成的两个同侧内角之和小于两直角，则延长这两条直线后，它们在有限距离处相遇。\cite[Book I, Axioms, XII]{Casey1885Euclid}这里的假设涉及第三条直线、内角以及角和比较；不能把所读条目的编号改写成“该版第五公设”。另一方面，Hilbert 的 Townsend 英译本以“过直线外一点恰有一条与该直线不相交的直线”表述平行公理，并随即把它拆为存在与唯一两个断言。\cite[Chapter I, §5]{Hilbert1950Foundations}

\ProseParagraphBreak

这两种表述提出的判定方式不同。前者由角的条件推出相交，后者直接限制不相交直线的条数。要在一个抽象体系内把它们互相推出，须先确定共同的其他公理，并检查角的作出、角与顺序的关系以及所用三角形结论怎样得到；把一句话换成另一句话本身不能代替这些推导。在实欧氏模型中，两种条件均可用已有坐标和角度计算检验，这也只说明它们在该模型中成立，不能单凭这个事实得到任意公理模型中的等价性。

\ProseParagraphBreak

因此，重述历史证明时应保留两条可查的联系：一条从具体版本通向它实际采用的词义、假设与推理；另一条从现代陈述通向所选公理或模型中的证明。前者帮助我们避免把编者补充写成原著文字，后者帮助我们说明补入的平方根、连续性或唯一性究竟解决了哪一个问题。等边三角形的例子所揭示的关键，正是“构造得到的点满足要求”与“所需的点确实存在”这两项证明任务。

\begin{chaptersummary}
	公理规定对象之间应满足的关系，模型为对象与关系提供解释。关联条件不决定长度，直线上的顺序也不自动提供连续性；增加合同与连续性条件时，需要明确新增了哪些可用的构造和结论。在坐标模型中，实数的代数、顺序与完备性分别承担不同的存在性依据。

	\ProseParagraphBreak

	欧氏平面的直线和上半平面的双曲直线都满足本附录列出的关联公理，而唯一平行命题在两者中取不同真值。推理的保真性由此排除了从这组公理推出该命题或其否定的可能。独立性总是相对于指定公理系统的；模型存在所依赖的背景理论也必须保留，因而相对一致性并不是一个无条件的绝对保证。

	\ProseParagraphBreak

	给定实仿射空间以后，标架把点对应到实数坐标，直线和介于关系也随之得到代数表达；增加内积后，Gram 矩阵再记录长度与角度。改变标架会改变分量，却不改变原几何对象。仅保留关联的变换可以改变合同关系，这说明选择坐标与选择几何结构是两个不同步骤。

	\ProseParagraphBreak

	阅读历史证明时，应先恢复原文中的对象、构造要求和已用命题，再比较现代重述。等边三角形的坐标构造给出了明确交点，但它使用的实数运算并非原文的逐步翻译。原始证明、补足存在性依据后的重述和另一种现代证明，各有明确的内容，不能只凭结论相同便视为同一个论证。
\end{chaptersummary}

\BookExercises
以下练习分别考察有限模型、结构之间的区别、坐标数据的恢复及构造所需的数域。每题都应先明确正在保留哪些关系，再判断结论能否推出。

\begin{exercise}
	\label{g1:ex:four-point-incidence}
	令 \(X=\{a,b,c,d\}\)，把 \(X\) 的每个二元子集都当作一条线，关联为点属于子集。
	\begin{enumerate}
		\item 核验三个关联公理 \(\mathsf I\)，并证明该模型满足唯一平行线命题 \(\mathsf P\)。
		\item 证明：在这个模型中，无法定义一种严格介于关系，同时使“介于的三点互异且共线”和“任意线段都能向一端延长”成立。
	\end{enumerate}
\end{exercise}
\begin{solution}
	(1) 两个不同点恰好确定由它们组成的二元子集，所以连线存在且唯一；每条线正有两个点；任取三个不同点，它们不可能同在一个二元子集中，故不共线。三条关联公理都成立。固定线 \(\ell\) 与其外一点 \(p\)，补集 \(X\setminus\ell\) 恰有两个点并包含 \(p\)，所以它是一条过 \(p\) 且与 \(\ell\) 不交的线。任何与 \(\ell\) 不交的线都是这个二元补集，故唯一。

	(2) 取任意两个不同点 \(a,b\)。若线段能延长，就应有点 \(e\) 使 \(B(a,b,e)\) 成立。严格介于要求 \(a,b,e\) 互异且共线，但经过 \(a,b\) 的唯一线只有 \(a,b\) 两点，不可能含有第三点。因此关联公理加上唯一平行线，并没有保证通常的线内顺序与延长性质。
\end{solution}

\begin{exercise}
	\label{g1:ex:seven-point-incidence}
	在二元域 \(\mathbb F_2=\{0,1\}\) 上取向量空间 \(V=\mathbb F_2^3\)。以 \(V\setminus\{0\}\) 为点集；对每个非零线性函数 \(f:V\rightarrow\mathbb F_2\)，取 \(\ker f\setminus\{0\}\) 为一条线。
	\begin{enumerate}
		\item 证明共有七点、七线，每条线有三个点，且满足 \(\mathsf I\)。
		\item 证明任意两条不同的线恰好有一个公共点。说明 \(\mathsf P\) 在此模型中的失败与双曲模型中的失败有什么不同。
	\end{enumerate}
\end{exercise}
\begin{solution}
	(1) \(V\) 有 \(2^3=8\) 个向量，去掉零向量后有七点。对偶空间也有八个元素，故有七个非零线性函数。每个非零 \(f\) 的核是二维空间，含四个向量，所以对应的线有三个点。不同的非零函数给出不同的核：若两个函数核相同，它们都在核上取零，并在核外取唯一的非零标量 \(1\)，于是函数相同。这也证明共有七条线。

	两个不同的非零向量 \(u,v\) 线性无关，因为 \(\mathbb F_2\) 的非零标量只有 \(1\)。经过两点的线必须由一个在 \(\operatorname{span}(u,v)\) 上为零的非零线性函数给出。这些函数构成对偶空间中一个一维子空间，因而恰有一个非零元素，证明连线存在且唯一。每线已有三个点；而标准基 \(e_1,e_2,e_3\) 不共线，因为非零函数不能同时在一组基上取零。故 \(\mathsf I\) 成立。

	(2) 若 \(f,g\) 非零且不同，则它们在对偶空间中线性无关。线性映射 \(v\mapsto(f(v),g(v))\) 的秩为 \(2\)，其核 \(\ker f\cap\ker g\) 维数为 \(1\)，因此恰含一个非零向量。该向量就是两线唯一的公共点。

	任取一条线及其外一点，所有过该点的线都与原线相交，没有不交线。因此这里破坏的是平行线的存在性。双曲模型的既有反例则给出了同一外点处至少两条不交线，破坏的是唯一性。两者都否定 \(\mathsf P\)，却不是同一种几何情形。
\end{solution}

\begin{exercise}
	\label{g1:ex:taxicab-betweenness}
	在 \(\mathbb R^2\) 上保留通常的仿射直线，改用距离
	\[
		d_1\bigl((x_1,y_1),(x_2,y_2)\bigr)=|x_1-x_2|+|y_1-y_2|.
	\]
	\begin{enumerate}
		\item 核验 \(d_1\) 是距离，并为每条仿射直线构造一个保持距离的实参数。
		\item 取 \(A=(0,0),B=(1,0),C=(1,1)\)，证明距离可加，但三点不共线。
		\item 以距离相等定义线段合同，找出两条 \(d_1\)-合同却并非欧氏等长的线段。说明点线关联相同为何不能决定合同关系。
	\end{enumerate}
\end{exercise}
\begin{solution}
	(1) 非负性、对称性以及距离为零恰在两点相同时，均来自实数绝对值。把两个坐标各自的三角不等式相加，得到 \(d_1(P,R)\le d_1(P,Q)+d_1(Q,R)\)。对直线 \(a+\mathbb R u\)，其中 \(u=(u_1,u_2)\ne0\)，令 \(c=|u_1|+|u_2|>0\)，取
	\[
		\gamma(t)=a+\frac{t}{c}u.
	\]
	它是从 \(\mathbb R\) 到该线的双射，且 \(d_1(\gamma(s),\gamma(t))=|s-t|\)。因此正文列出的线内顺序、唯一复制长度与序完备性仍成立。

	(2) 三段距离为 \(d_1(A,B)=1\)、\(d_1(B,C)=1\)、\(d_1(A,C)=2\)，满足可加等式。然而 \(B-A=(1,0)\) 与 \(C-A=(1,1)\) 线性无关，所以三点不共线。若删除正文定义中独立列出的共线条件，就会错误地把 \(B\) 宣布为介于 \(A,C\) 的点。

	(3) 取从原点到 \((1,0)\) 及从原点到 \((1/2,1/2)\) 的两段。它们的 \(d_1\)-长度都是 \(1\)，而欧氏长度分别为 \(1\) 与 \(1/\sqrt2\)。两种距离是在同一仿射点线结构上增加的不同数据，所以保留全部关联关系仍可改变合同判断。
\end{solution}

\begin{exercise}
	\label{g1:ex:model-inference-limits}
	设 \(T\) 是使用可靠的一阶推演规则的公理系统，\(\varphi\) 是同语言中的句子，模型 \(\mathcal M,\mathcal N\) 满足
	\[
		\mathcal M\models T\cup\{\varphi\},\qquad
		\mathcal N\models T\cup\{\neg\varphi\}.
	\]
	\begin{enumerate}
		\item 证明 \(\varphi\) 对 \(T\) 独立，并说明两个扩充系统的一致性依据。
		\item 有人说：“因为 \(T\) 有模型，所以 \(T\) 的所有模型都同构。”以上材料能否支持这一结论？
		\item 设 \(T'=T\cup\{\varphi\}\)。解释为何 \(\varphi\) 对 \(T\) 独立，却不对 \(T'\) 独立；指出原来的模型反证在哪一步不能用于 \(T'\)。
	\end{enumerate}
\end{exercise}
\begin{solution}
	(1) 若 \(T\vdash\varphi\)，保真性使 \(\mathcal N\models\varphi\)，与 \(\mathcal N\models\neg\varphi\) 矛盾；若 \(T\vdash\neg\varphi\)，同理与 \(\mathcal M\models\varphi\) 矛盾，所以两个方向都不能证明。每个扩充系统都有题设给出的模型；若其中一个推出某句子及其否定，该模型就会同时满足二者。因此在构造模型并验证推演保真性的背景理论中，两个系统均一致。

	(2) 不能。对同一种语言，同构在每个基本关系上保持真假，因而对由这些关系经逻辑连接和量化组成的句子也保持真假。具体地，连接词的步骤由真值表保证；存在量词的见证可由同构映到另一模型，全称量词则由同构的满射性保证。对句子的构成归纳，便知同构模型必须在 \(\varphi\) 上同真同假。题设的 \(\mathcal M,\mathcal N\) 却在 \(\varphi\) 上取相反真值，所以它们恰好给出了不可能同构的两个模型。

	(3) \(\varphi\) 已是 \(T'\) 的一条公理，故有一行证明 \(T'\vdash\varphi\)，独立性失效。\(\mathcal N\) 虽然满足 \(T\)，却不满足新加入的 \(\varphi\)，因此不是 \(T'\) 的模型。把它用于排除从 \(T'\) 出发的证明，违反了保真性命题中的模型假设。这说明独立性结论必须连同基准公理系统一起陈述。
\end{solution}

\begin{exercise}\label{g1:ex:oblique-frame-model}
	在通常欧氏平面中取原点 \(o=(1,-1)\)，方向基 \(e_1=(2,0)\)、\(e_2=(1,1)\)，记一个点的新坐标为 \((u,v)\)。求此标架中的内积矩阵、点 \(p=(5,1)\) 的坐标及 \(d(o,p)\)，并用 \((u,v)\) 写出线段 \(op\) 的中垂线方程。若方向基改为 \(e'_1=e_1\)、\(e'_2=e_1+e_2\)，原点不变，求新内积矩阵及 \(p\) 的新坐标，并核对长度不变。
\end{exercise}
\begin{solution}
	原来的直角坐标与 \((u,v)\) 的关系为
	\[
		(x,y)=(1,-1)+u(2,0)+v(1,1)=(1+2u+v,-1+v).
	\]
	所以 \(p\) 的坐标为 \((1,2)\)。基向量的配对给出
	\[
		G=\begin{pmatrix}4&2\\2&2\end{pmatrix},\qquad
		d(o,p)^2=\begin{pmatrix}1&2\end{pmatrix}G\begin{pmatrix}1\\2\end{pmatrix}=20.
	\]
	因此 \(d(o,p)=2\sqrt5\)。若把这套非正交坐标误当成标准正交坐标，会错误地得到 \(\sqrt{1^2+2^2}=\sqrt5\)。

	令 \(x=(u,v)^{\mathsf T}\)、\(a=(1,2)^{\mathsf T}\)。与 \(o,p\) 等距的条件为 \(x^{\mathsf T}Gx=(x-a)^{\mathsf T}G(x-a)\)，消去二次项得到 \(2x^{\mathsf T}Ga=a^{\mathsf T}Ga\)。由于 \(Ga=(8,6)^{\mathsf T}\)，中垂线为 \(4u+3v=5\)。此处的线性方程来自度量比较，而不是把纸面坐标轴预设为垂直。

	换基矩阵为 \(S=\begin{psmallmatrix}1&1\\0&1\end{psmallmatrix}\)，满足旧坐标列等于 \(S\) 乘新坐标列。因此
	\[
		G'=S^{\mathsf T}GS=\begin{pmatrix}4&6\\6&10\end{pmatrix},
		\qquad p'=S^{-1}\begin{pmatrix}1\\2\end{pmatrix}=\begin{pmatrix}-1\\2\end{pmatrix}.
	\]
	直接计算 \((p')^{\mathsf T}G'p'=4-24+40=20\)，恢复同一个长度平方。
\end{solution}

\begin{exercise}\label{g1:ex:congruence-determines-scale}
	设 \(b_1,b_2\) 是同一个非零有限维实向量空间 \(V\) 上的两个正定内积，令 \(q_j(v)=b_j(v,v)\)。假设它们给出的等长关系相同，即对全部 \(u,v\in V\)，有
	\[
		q_1(u)=q_1(v)\quad\Longleftrightarrow\quad q_2(u)=q_2(v).
	\]
	证明存在常数 \(c>0\) 使 \(b_2=cb_1\)。若再指定一个非零向量在两个度量下长度相同，证明两个内积必相同。解释这与“合同关系还没有指定长度单位”的联系。
\end{exercise}
\begin{solution}
	选 \(q_1(e)=1\) 的向量 \(e\)，令 \(c=q_2(e)>0\)。任意满足 \(q_1(u)=1\) 的向量都与 \(e\) 对 \(q_1\) 等长，故由假设也有 \(q_2(u)=c\)。对非零 \(v\)，将其写成 \(v=\sqrt{q_1(v)}\,u\)，其中 \(q_1(u)=1\)，利用二次齐次性得到
	\[
		q_2(v)=q_1(v)q_2(u)=c\,q_1(v).
	\]
	零向量也满足此式。由极化恒等式，
	\[
		b_2(u,v)=\frac{q_2(u+v)-q_2(u)-q_2(v)}2
		=c\,b_1(u,v).
	\]
	若某个 \(v_0\ne0\) 在两种度量下长度相同，则 \(q_2(v_0)=q_1(v_0)>0\)，代入上式得 \(c=1\)。在已有实仿射和内积结构的范围内，全部线段合同关系因而确定内积到一个正比例因子；再用一条非零线段校准单位，才确定具体的数值长度。
\end{solution}

\begin{exercise}\label{g1:ex:frame-distances-rigidity}
	设 \(f:A\to B\) 是两个 \(n\) 维实欧氏仿射空间之间的连续中点保持映射，其中 \(n\ge1\)。取 \(A\) 的仿射标架 \((p_0,\ldots,p_n)\)，并假设
	\[
		d_B(f(p_i),f(p_j))=d_A(p_i,p_j)\qquad(0\le i,j\le n).
	\]
	证明 \(f\) 保持任意两点的距离。若只要求 \(p_0p_i\)、\(1\le i\le n\)，这 \(n\) 条线段的长度保持，结论是否仍成立？给出二维反例。
\end{exercise}
\begin{solution}
	由\cref{g1:prop:continuous-midpoint-affine}，\(f\) 有实线性部分 \(L\)。令 \(e_i=p_i-p_0\)，这些向量组成 \(A\) 的方向空间的基。给定距离等式先给出 \(\lVert Le_i\rVert=\lVert e_i\rVert\)，又给出 \(\lVert Le_i-Le_j\rVert=\lVert e_i-e_j\rVert\)。用极化恒等式得到
	\[
		\langle Le_i,Le_j\rangle_B
		=\frac{\lVert Le_i\rVert_B^2+\lVert Le_j\rVert_B^2-\lVert Le_i-Le_j\rVert_B^2}{2}
		=\langle e_i,e_j\rangle_A.
	\]
	对基展开并用双线性，可知 \(L\) 保持所有向量的内积，特别保持长度。又因 \(f(q)-f(p)=L(q-p)\)，任意两点距离均被保持。正定性还说明 \(L\) 单射，同维有限维性使它可逆，所以 \(f\) 实际为等距同构。

	只保留从 \(p_0\) 出发的边长不足以控制不同基方向之间的夹角。在通常欧氏平面中取标准标架 \(p_0=(0,0),p_1=(1,0),p_2=(0,1)\)，并令
	\[
		f(x,y)=\bigl(x+\tfrac12y,\tfrac{\sqrt3}{2}y\bigr).
	\]
	它线性，因而连续且保持中点。两向量 \(f(p_1)=(1,0)\)、\(f(p_2)=(1/2,\sqrt3/2)\) 都长为 \(1\)，满足所述 \(p_0p_i\) 的边长条件；但它们之间的距离为 \(1\)，而 \(d(p_1,p_2)=\sqrt2\)。所以 \(f\) 不是等距映射，遗漏的标架边恰好承载了夹角信息。
\end{solution}

\begin{exercise}\label{g1:ex:one-circle-construction-not-all}
	令 \(K=\mathbb Q(\sqrt3)\subset\mathbb R\)，以 \(K^2\) 为点集，距离平方使用通常的 \((x_1-y_1)^2+(x_2-y_2)^2\)。证明每一条端点在 \(K^2\) 中的非零线段都能在 \(K^2\) 中作为某个等边三角形的一边。再证明以 \((0,0)\)、\((1,0)\) 为圆心、共同半径为 \(\sqrt3\) 的两个圆周，在 \(K^2\) 中各自非空却没有共同点。解释为什么完成第一项构造仍不足以保证所有实平面中相交的圆在这个点集中相交。
\end{exercise}
\begin{solution}
	映射 \(J(x,y)=(-y,x)\) 保持 \(K^2\)，而 \(1/2\)、\(\sqrt3/2\) 都属于 \(K\)。因此对 \(A,B\in K^2\)、\(A\ne B\)，式\eqref{g1:eq:equilateral-coordinate-intersections}给出的两个点仍在 \(K^2\) 内，且同样的距离平方计算给出三边相等。非共线性也由垂直分量非零得到。

	第二组圆周各含点 \((\sqrt3,0)\) 与 \((1+\sqrt3,0)\)，故分别非空。共同点必须满足
	\[
		x^2+y^2=3,\qquad (x-1)^2+y^2=3.
	\]
	两式相减给出 \(x=1/2\)，代回得 \(y^2=11/4\)。如果 \(y\in K\)，则 \(\sqrt{11}\in K\)。设 \(\sqrt{11}=a+b\sqrt3\)、\(a,b\in\mathbb Q\)，平方后由 \(1,\sqrt3\) 在 \(\mathbb Q\) 上线性无关得到 \(2ab=0\) 和 \(a^2+3b^2=11\)。若 \(b=0\)，则 \(a^2=11\)；若 \(a=0\)，则 \(b^2=11/3\)。这两种情形都不可能，因为一个非零有理数的平方在约分后的分子、分母中，每个素因子的指数都为偶数，而 \(11\) 和 \(11/3\) 不具备该性质。因此两圆周没有 \(K^2\) 中的共同点。

	在 \(\mathbb R^2\) 中，这两个圆确实交于 \((1/2,\pm\sqrt{11}/2)\)。第一项构造只需要 \(\sqrt3\)，第二项交点却需要额外的平方根 \(\sqrt{11}\)。一种作图能够完成，并不意味着其他作图所需的代数运算也已经在同一个点集中封闭。本题没有断言 \(K^2\) 满足一套完整的中性几何公理。
\end{solution}

